\documentclass[a4paper,12pt]{article}

\usepackage{cmbright}
\usepackage[latin1]{inputenc}
\usepackage[T1]{fontenc}
\usepackage[a4paper]{geometry}
\usepackage[dvips]{graphicx}
\usepackage[dvips]{graphics}
\usepackage{times}
\usepackage{latexsym}
\usepackage{url}


\renewcommand{\arraystretch}{1.2}
\pagestyle{empty}


\newcommand{\ligne}[1]{\rule[0.5ex]{\textwidth}{#1}\\}
\newcommand{\activite}[1]{\textit{#1}\ }
  
\newcommand{\interRubrique}{\bigskip}
\newcommand{\styleRub}[1]{\textbf{\large #1}\par}
\newcommand{\indentStd}{\noindent\hspace*{10pt}}



\newenvironment{rubrique}[2][\linewidth]% "rubrique" prend deux arguments
{\styleRub{#2}%       			% le second argument : titre
\ligne{0.5mm}
\setlength{\lenB}{#1}%		       	% Le premier argument :indentation
\setlength{\lenC}{\linewidth}%		% Calculs...
\addtolength{\lenC}{-\lenA}%
\addtolength{\lenC}{-\lenB}%
%\addtolength{\lenC}{-\parindent}%
\addtolength{\lenC}{-19pt}
\indentStd\begin{tabular}[t]{p{\lenB}p{\lenC}}}
{\end{tabular}}

\newenvironment{rubriquesansligne}[2][\linewidth]% "rubrique" prend deux arguments
{\styleRub{#2}%       			% le second argument : titre
\setlength{\lenB}{#1}%		       	% Le premier argument :indentation
\setlength{\lenC}{\linewidth}%		% Calculs...
\addtolength{\lenC}{-\lenA}%
\addtolength{\lenC}{-\lenB}%
\addtolength{\lenC}{-\parindent}%
\addtolength{\lenC}{-19pt}
\indentStd\begin{tabular}[t]{p{\lenB}p{\lenC}}}
{\end{tabular}}

\newlength{\lenA} % indentation au début d'une ligne
\setlength{\lenA}{0.cm}
\newlength{\lenB} % Taille champ dates
\newlength{\lenC} % Taille champ description
%\newlength{\lenD}


\geometry{hmargin={2cm, 2cm}, vmargin={2cm, 2cm}}

\begin{document}


\begin{tabular*}{1\textwidth}{@{\extracolsep{\fill}}lr}
Mathieu Lacage & mathieu.lacage@cutebugs.net\\
\end{tabular*}

\vspace{1cm}
\begin{center}{\huge Software Engineer}\end{center}

\begin{center}
Experienced Software Engineer who enjoys hands-on technical roles.\newline
I build secure stuff that works and that scales.
\end{center}

\vspace{1cm}
\begin{rubrique}{Technical Experience}

Technology
\begin{itemize}
\item Languages: C, C++, Python, SQL
\item Infrastructure tools: MySQL, Ansible, uWSGI, Prometheus, Grafana
\end{itemize}

Extensive practical and theoretical knowledge of
\begin{itemize}
\item Networking: from Wi-Fi MAC, to IP and TCP
\item Machine Learning, Transformers
\item Event-driven discrete-time simulation
\end{itemize}

\end{rubrique}

\vspace{2cm}
\begin{rubrique}{INRIA, 2025 -- Present, Research Engineer}

Researching domain-specific SLMs (Small Language Models) under 1B parameters

\begin{itemize}
\item Dataset selection and tokenizer adaptation
\item Quantization and its impact on training convergence and speed
\item Training vs inference cost tradeoffs and their impact on model architecture
\end{itemize}

Select software projects
\begin{itemize}
\item Interactive visualization of cost, pricing, and margin models for LLM inference
\item Container-based single-GPU SLM training framework for students
\end{itemize}

\end{rubrique}
\pagebreak

\vspace{0.5cm}
\begin{rubrique}{Hive, 2023 -- 2025, Architect}

Scaled (100x) and improved the reliability (2x) of a Go libp2p-based application.

\begin{itemize}
\item Metric collection, analysis, monitoring, and alerting
        spanning centralized infrastructure and a p2p network
\item Debugged and fixed network and IO problems to scale the Hive p2p network
\item Built python-based development, testing, and debugging tools
\item Planned and managed component-level testing strategies
\item Designed forward error correction policies around dynamic
        decisions to minimize both data loss and storage use
\end{itemize}

\end{rubrique}

\vspace{0.5cm}
\begin{rubrique}{Alcmeon, 2011 -- 2023, CTO, co-founder}

As CTO, I 
\begin{itemize}
\item Made sure our customers's data stayed safe via at-rest encryption,
        and centralized access control
\item Scaled our code and infrastructure to sustain
        50\% YoY volume growth with 1000\% YoY peak volume growth
\item Recruited and led the technical team (8 years)
\item Mentored and directly managed the engineering manager and product
        director
\end{itemize}


As co-founder, I focused on
\begin{itemize}
\item Delivering high-quality customer support
\item Product design to ensure we build features that grow our business
\item Presales to help close deals with prospects
\end{itemize}

Select projects
\begin{itemize}
\item Pub/Sub streaming WebSocket server that scales to 10K concurrent connections per cpu
\item Zero-downtime scale up of a single database with a single table to partioned
 databases per customer, 250 tables, 1TB
\item SSO integration in SSH
\item A security policy that spans infrastructure and application-level code
\end{itemize}

\end{rubrique}
\pagebreak


\begin{rubrique}{INRIA, 2005 -- 2011, Software Lead ns-3}
  Built ns-3, an open-source network simulator now used in hundreds of research publications every year. 
 
  \begin{itemize}
  \item Designed and implemented core APIs: object model, network packets, event scheduler
  \item Implemented models for UDP/IP/ICMP, MAC/PHY Wi-Fi network protocols
  \item Integrated unmodified real-world network protocol implementations
  \end{itemize}

  But also:
  \begin{itemize}
  \item Recruited, and managed a local development team (5).
  \item Relocated to the University of Washington for 10 months to initiate collaboration with US team.
  \item Evangelized use of ns-3 within other research institutes via
    presentations and seminars.
  \item Published 7 papers on ns-3 (4 as main author and 3 as co-author)
  \end{itemize}

\end{rubrique}

\vspace{0.5cm}
\begin{rubrique}{INRIA, 2003 -- 2005, Software Engineer}

  Designed and built software for network research teams:  
  \begin{itemize}
  \item Yans, a C++ event-driven simulator
  \item NEPI, a python tool used to describe, deploy, and control networking
    experiments on hundreds of hosts distributed all over the planet.
  \end{itemize}

  Provided software mentoring to research projects involved in 
  bio-reactor chemical reaction control and medical image analysis.

\end{rubrique}

\vspace{0.5cm}
\begin{rubrique}{Sigma-Designs, 2001 -- 2003, Software Engineer}

  Implemented DVD navigation control software for the video
  decompression chips that were designed in-house and
  sold to OEMs to build consumer DVD players.

\end{rubrique}

\vspace{0.5cm}
\begin{rubrique} [4cm] {Education}
  2006 -- 2010 & Ph.D. at University of Nice,
  \emph{Experimentation Tools for Networking Research}, under the supervision of Walid Dabbous\\
  1998 -- 2001 & Engineer at Telecom ParisTech (ENST), Software Engineering, Networking, Micro-Electronics \\
\end{rubrique}

\end{document}


